% status: ZERO
% Chapter 9 of the second edition. Plan: docs/STRUCTURE.md. Rules: AUTHORING.md.
% Measurements: research/measurements/2026-09-23-m1-part2.md
\chapter{Decode Attention: FlashDecoding, Split-K and Paged Kernels}\label{ch:decode-attention}

\begin{ChapterIntent}
At decode each sequence contributes a single query, so a prefill-style
attention kernel leaves most of the GPU idle while one block walks a long
cache. This chapter shows how decode attention restores parallelism by
splitting the key/value sequence and merging partial softmaxes through their
log-sum-exp, how grouped-query attention changes what can be shared, and what
a paged cache costs the kernel. Its central lesson is measured: a
decode-attention kernel can be limited by occupancy rather than bandwidth, and
then an optimization that saves bytes by sacrificing parallelism makes it
slower.
\end{ChapterIntent}

\OpeningSection{Twice the bytes, twice the time}

On the M1 of \Chref{roofline}, llama.cpp decodes Llama~3.2~3B at 25.0 tokens
per second with an empty context. With 16,384 tokens already in the cache it
decodes at \textbf{12.4} (measured with flash attention enabled; record
\texttt{2026-09-23-m1-part2}). The extra work is attention over the cache, and
its bytes are easy to count: 16,384 tokens at 112~KiB is 1.88~GB of keys and
values --- nearly as much as the model's 2.01~GB of weights, and read at every
step beside them. The step reads 1.93 times as many bytes and takes 2.03 times
as long.

At every depth measured --- 0, 4,096, 8,192 and 16,384 tokens --- the step time
matches its bytes read at 48--52~GB/s: the rate at which the same engine reads
its weights, and about what a plain streaming kernel reaches on this GPU. A
long cache costs what its bytes cost, and no more. That is an achievement, not
a law. Each step asks the kernel to read 1.88~GB on behalf of a single query
vector per head, and the natural way to write that kernel leaves most of the
GPU idle. Another engine's kernel on the same M1 read its cache at about a
seventh of this rate (the section on occupancy, below). This chapter is about
the difference.

(Absolute rates on this laptop move with its load. A first run of these
measurements, taken while the machine was heavily loaded, recorded lower rates
that fell further at 16,384 tokens, and \Chref{roofline}'s 21 tokens per
second came from a busier session than this one. The record keeps every run
with its conditions; this chapter quotes the quietest.)

\NapkinSection{Napkin math: parallelism in one query}

For one sequence, one layer and one key/value head, decode attention reads
$S$ cached keys and values, $2Sd_hb_{kv}$ bytes, and performs about
$4Sd_h$ FLOPs for each of the $g=H_q/H_{kv}$ query heads that share them. Its
arithmetic intensity is therefore
\begin{equation}
I_{\text{attn}} \approx \frac{4\,g\,S\,d_h}{2\,S\,d_h\,b_{kv}} = \frac{2g}{b_{kv}}
\quad[\text{FLOP/byte}],
\label{eq:attn-intensity}
\end{equation}
one FLOP per byte for multi-head attention in 16-bit, four for a group of four,
independent of context length. Decode attention is a bandwidth-bound workload,
and batching does not change that: every sequence reads its own cache.

The parallelism is the other number. A prefill kernel spreads work over query
tokens; at decode there is one. The natural decomposition --- one thread block
per (sequence, head) --- gives $B\,H_q$ blocks. For Llama~3.2~3B at batch one
that is 24 blocks, against a GPU that needs hundreds of concurrently resident
warps to hide memory latency. Each block then walks the entire cache alone.
The sequence dimension is the only place left to find parallelism.

\section{The shape of decode attention}

For one head, decode attention is two matrix--vector products around a softmax:
scores $\mathbf{s}=\mathbf{K}\mathbf{q}/\sqrt{d_h}$ with
$\mathbf{K}\in\mathbb{R}^{S\times d_h}$, probabilities
$\mathbf{p}=\operatorname{softmax}(\mathbf{s})$, output
$\mathbf{o}=\mathbf{V}^{\top}\mathbf{p}$. No score matrix exists to be
avoided --- the scores are a single vector of length $S$ --- so the problem
\Chref{flashattention} solved is not this one. The problem here is reading
$\mathbf{K}$ and $\mathbf{V}$ as fast as memory allows with far too few
independent tasks.

\section{Splitting the sequence}

Flash-Decoding divides the cache into chunks and gives each chunk to its own
thread block \cite{dao2023flashdecoding}. Each block computes attention over its
chunk exactly as \Chref{flashattention}'s online softmax would, producing a
partial output with its running maximum and denominator. A second, small pass
merges the partials. With 16 chunks, a batch-one decode step has 16 times the
parallelism. Its authors reported attention up to 50 times faster than
FlashAttention at very long contexts and end-to-end decoding up to 8 times
faster at batch one on CodeLlama-34B (their measurements, 2023).

\section{Merging partial softmaxes}

Two partials over disjoint chunks, each with maximum $m_i$, denominator
$\ell_i$ and unnormalized output $\mathbf{o}_i$, combine exactly:
\begin{equation}
m = \max(m_1,m_2),\quad
\ell = \ell_1e^{m_1-m}+\ell_2e^{m_2-m},\quad
\mathbf{o} = \mathbf{o}_1e^{m_1-m}+\mathbf{o}_2e^{m_2-m},
\label{eq:lse-merge}
\end{equation}
and the result is $\mathbf{o}/\ell$. This is equation~\eqref{eq:online-softmax}
applied to whole chunks instead of single scores; engines often store
$m+\log\ell$, the log-sum-exp, as a single number per partial. The merge is
associative, so partials can be combined in any grouping --- but floating-point
rounding makes the \emph{order} observable, and the order becomes part of the
engine's reproducibility contract (\Chref{fast-wrong-answers}).

\section{Grouped-query attention: what the heads share}

With grouped-query attention, the $g$ query heads of a group read the same keys
and values. A kernel can process the whole group in one pass over the cache,
reading each key once instead of $g$ times, and equation~\eqref{eq:attn-intensity}
says that multiplies the useful work per byte by $g$. On a CPU core, which has a
fixed amount of parallelism and pays for every cache miss, packing the group
this way is a clear gain. On a GPU the same change also divides the number of
thread blocks by $g$, and whether it helps depends on whether bytes or
parallelism was the limit. The case study of this chapter is what happened when
that question was answered by experiment.

\section{Reading through a block table}

A paged cache (\Chref{paged-kv}) stores each sequence's keys and values in
fixed-size blocks scattered through a pool, found through a per-sequence table.
The kernel's inner loop gains an indirection: for each block of tokens, look up
where it lives. The table is small and the lookups are cheap compared with the
bytes they locate, provided blocks are large enough to be read in efficient
bursts. Attention libraries such as FlashInfer represent paged, ragged and
sparse layouts in one block-sparse format so that one kernel family serves them
all \cite{ye2025flashinfer}; DeepSeek's FlashMLA serves paged latent attention
\cite{deepseek2025flashmla}.

\section{Occupancy, not bandwidth}

A GPU hides memory latency by switching among many resident warps; a kernel
that launches too few blocks cannot keep enough requests in flight to use the
memory's bandwidth, however few bytes it needs. That is the regime of decode
attention at small batch and moderate context, and it inverts the usual
optimization rule. Reducing bytes helps a kernel that is waiting on bandwidth.
It does nothing for a kernel that is waiting for work --- and if reducing bytes
also reduces the number of blocks, it makes the kernel slower.

The same GPU shows both kinds of kernel. llama.cpp's decode attention reads the
cache at the 48--52~GB/s of the opening measurement, the bandwidth roof in
practice. Hermon's Metal kernel, measured on the decode shape of Llama-3-8B,
took 19.2~ms to read one layer's 134~MB of 16-bit keys and values at 32,768
tokens: about 7~GB/s (derived from Hermon \texttt{docs/KERNEL\_DESIGN.md},
section K4). Both read the same bytes. What differs is the kernel, and the
case study below is its developers' account of what they learned about it.

\section{Fixed split count or fixed split size}

How many chunks should a sequence be split into? A fixed \emph{count} gives
every sequence the same parallelism, but the chunk boundaries --- and so the
floating-point order of the merge --- move with the length. A fixed
\emph{size}, say 256 tokens per chunk, keeps the boundaries in place: the
split, and so the result, depends only on the tokens. Thinking Machines showed
that fixed-size splits are one of the changes needed to make attention's
result independent of which other requests share the batch
\cite{he2025nondeterminism}; Hermon uses a fixed size of 256 tokens, clamped to
between 1 and 64 splits, so that its plans depend only on the shape of the
request (Hermon \texttt{docs/KERNEL\_DESIGN.md}, section 6.3).

\LedgerSection

Splitting decode attention along the sequence and merging the partials:

\begin{ConstraintLedger}
Latency & buys & Parallelism for batch-one decode: 16 splits give 16 times as many independent blocks. \\
Memory bandwidth & moves & The same cache bytes, now read fast enough to approach the bandwidth roof. \\
Memory capacity & spends & A workspace of partial outputs, maxima and denominators per split. \\
Compute & spends & A merge pass per head per step. \\
Correctness & risks & The merge order is observable in the last bits; fix it by construction. \\
\end{ConstraintLedger}

\LabSection{split-K decode attention with a log-sum-exp merge}

\LabIntent{In \texttt{code/labs/ch09-decode-attention/}, a Rust decode-attention
kernel over a synthetic cache, first single-threaded, then split into $k$
chunks run on a thread pool and merged with equation~\eqref{eq:lse-merge}. A
float64 oracle checks every configuration. The reader measures throughput as
$k$ grows, then implements grouped-query packing and finds whether it helps
on their CPU --- and predicts, before reading the case study, whether it would
help on a GPU.}

\HappenedSection{the GQA fusion that lost in every shape}

Hermon's attention kernel for Apple GPUs gave each query head its own thread
group. Under grouped-query attention that meant the heads of a group each read
the same keys and values --- four times over at 32 query heads and 8
key/value heads. On the CPU, packing a group to read the cache once was a real
win, so the GPU analogue looked certain to help. Hermon's developers built the
fused kernel, verified that it was bit-identical to the per-head kernel, and
measured it on an M1. It was \textbf{slower in every shape measured}:
0.19--1.00 times the per-head kernel's speed at decode across four head
configurations, and 0.68--0.71 times at prefill (Hermon
\texttt{docs/KERNEL\_DESIGN.md}, section K4.1).

Both reasons invalidated the premise. Fusing divided the number of thread
groups by the group size, and at decode the grid was already too small to fill
the device. And the ``redundant'' reads had never cost bandwidth: the thread
groups of a group ran concurrently and read the same cache lines, so the
hardware had already removed the traffic the fusion set out to save. The
developers' conclusion generalizes: on that device, decode attention is
occupancy-bound, not bandwidth-bound, and an optimization that trades
parallelism for bytes will lose. They recorded the failure so that it would not
be repeated --- which is the only reason we can learn from it.

\FrontierSection{2026-09-24}

Decode attention kernels are now chosen per hardware and per model family.
vLLM's documented defaults put FlashAttention first on Hopper and FlashInfer
first on Blackwell, with separate backends for latent attention and for
DeepSeek-V4's sparse attention \cite{vllm2026attentionbackends}; FlashMLA's
dense decode kernel reports up to 3,000~GB/s on an H800
\cite{deepseek2025flashmla}. Sparse decode attention, which reads only selected
parts of the cache, changes this chapter's arithmetic (\Chref{long-context}), and
attention over quantized caches changes the bytes (\Chref{kv-quantization}).

\ExercisesSection

\begin{enumerate}
\item \emph{Derivation.} Prove that equation~\eqref{eq:lse-merge} is
associative, and give a floating-point example where two groupings produce
different results.
\item \emph{Prediction.} For Llama-3-8B at batch one and 4,096 tokens of
context, how many 256-token splits does each head get, and how many
independent blocks does the step launch?
\item \emph{Implementation.} In the lab, measure the workspace memory of the
split partials as a function of batch, heads and splits.
\item \emph{Falsification.} Find a context length and thread count on your CPU
at which splitting makes decode attention slower.
\end{enumerate}
